% !PRINTED 463
\weberpage{463}
%<<CONT
$17^{\text{ten}}$ Grades, die nicht existirt, und folglich ist $K'$ mit
$aK$ identisch.

\begin{center}
\gesperrt{Es giebt also, von einem constanten Factor abgesehen, nur
eine Invariante $21^{\text{sten}}$ Grades.}
\end{center}

Nun ist aber das System der 21 Hauptaxen auch eine Invariante
$21^{\text{sten}}$ Grades, und daraus ist zu schliessen:

\begin{center}
\gesperrt{Die Invariante $K$ zerfällt in 21 lineare Factoren, die,
gleich Null gesetzt, die Hauptaxen der Gruppe darstellen.}
\end{center}

Wir können sodann eine Invariante $42^{\text{sten}}$ Grades bilden,
nämlich
\weq{2}{\Delta^7 - kC^3 = Q,}
worin $k$ eine beliebige Constante ist, und diese Constante lässt
sich so bestimmen, dass die Curve $Q$ durch einen beliebig gegebenen
Punkt auf $f$ geht, und sie muss dann auch durch alle mit diesem
verbundenen Punkte hindurchgehen. Also können wir $k$ in $Q$ so
bestimmen, dass die Curve $Q$ aus der Curve $f$ ein beliebig
gegebenes System verbundener Punkte ausschneidet.

Ist nun $\Phi$ eine beliebige durch $f$ nicht theilbare Invariante,
die $h_1, h_2, h_3$ mal durch die Pole $P_7, P_3, \overline{P}_2$
geht, und ausserdem durch beliebige Systeme $S_1, S_2, \ldots$
verbundener Punkte auf $f$, die auch theilweise zusammenfallen
können, so bilden wir zunächst nach (2) die Formen
$Q_1, Q_2, \ldots$, die in den Systemen $S_1, S_2, \ldots$
verschwinden, und dann die Form
\[
\Psi = \Phi - a\Delta^{h_1} C^{h_2} K^{h_3} Q_1 Q_2 \ldots
\]

Die Curve $\Psi$ geht für jeden Werth der Constanten $a$ durch die
sämmtlichen Schnittpunkte von $\Phi$ mit $f$, und wenn wir also $a$ so
bestimmen, dass $\Psi$ durch irgend einen davon verschiedenen Punkt
von $f$ geht, so ist $\Psi$ durch $f$ theilbar, also
\[
\Phi = a\Delta^{h_1} C^{h_2} K^{h_3} Q_1 Q_2 \cdots + f\Phi_1.
\]
Darin ist nun $\Phi_1$ wieder eine Invariante, aber von niedrigerem
Grade als $\Phi$, und durch vollständige Induction ist hiermit der
Satz bewiesen:

\begin{center}
\gesperrt{Jede Invariante der Gruppe lässt sich als ganze rationale
Function der vier fundamentalen Invarianten $f, \Delta, C, K$
darstellen.}
\end{center}

Bestimmt man in (2) die Constante $k$ so, dass die Curve $Q$ durch
einen der Pole $\overline{P}_2$ geht, so kann sie durch keinen nicht
%CONT>>
